% ============================================================
% ［架空の記入例・提出不可］研究設定，症例数，数値および結果はすべて架空である．
% 引用文献は実在するが，本例の研究成果を裏づけるものではない．
% 通常のchapters/01-introduction.texには自動で読み込まれない．
% ============================================================
\chapter{\UsamiIntroductionTitle}

\begin{center}
  \UsamiFictionalExampleNotice
\end{center}

\section{研究背景}

腹部CTにおける多臓器領域分割は，臓器体積の計測，治療計画および画像解析の前処理に用いられる．U-Net\cite{ronneberger2015unet}は画像の局所情報と意味情報を統合し，3D U-Net\cite{cicek20163dunet}はこの考え方を体積データへ拡張した．一方，施設ごとの撮像条件，造影相およびアノテーション方針の違いにより，十分な症例を同一条件で収集することは容易でない．特に小臓器や細い構造では，症例不足が分割結果の不安定さとして現れる可能性がある．

幾何変換や濃度変換によるデータ拡張は，画像とラベルへ同じ変換を適用できる反面，実データにない形状や濃度分布を新たに生成する能力は限られる．Denoising Diffusion Probabilistic Models\cite{ho2020ddpm}および潜在拡散モデル\cite{rombach2022ldm}は高品質な画像生成を可能にするが，医用画像へ適用する際には，視覚的に自然な画像が必ずしも臓器の解剖学的妥当性を満たすとは限らない．また，条件ラベルと生成画像の境界が対応しなければ，拡張データが分割器へ誤った教師情報を与える．

\section{研究目的}

本研究の目的は，少数症例下の腹部CT多臓器分割を対象として，画像と臓器ラベルの対応を保つ条件付き潜在拡散データ拡張を構成し，解剖構造の検査を通過した合成対だけを分割学習へ利用する方法を検討することである．臓器体積，左右臓器の位置関係および臓器間重複を検査し，実画像のみ，幾何拡張，検査を伴わない拡散拡張と患者単位で比較する．評価では臓器平均Dice係数と95パーセンタイルHausdorff距離（HD95）を併用し，重複精度と境界誤差を分けて確認する．

\section{本研究の貢献}

本例で想定する貢献は3点である．第1に，腹部CTと多臓器ラベルを条件とする潜在拡散データ拡張へ，臓器間関係を検査する採択処理を導入した．第2に，実画像と採択合成画像を用いる分割器へ，解剖学的に不自然な予測を抑制する整合性項を組み込んだ．第3に，患者単位のデータ分割，DiceとHD95，除去実験および失敗例を含む評価設計を示した．

これらの貢献，症例設定および性能は，修士論文の書き方を示すための架空例であり，臨床的有効性を主張するものではない．

\section{本論文の構成}

第2章では多臓器分割と拡散モデルによるデータ拡張を整理し，第3章で対象データ，生成対の採択条件および評価課題を定義する．第4章では解剖構造保持拡散拡張と分割学習の目的関数を示す．第5章で架空の比較実験と失敗例を報告し，第6章で得られた知見，限界および今後の課題を述べる．
